%! Author = sriadityavedantam
%! Date = 3/30/26

\section{Constructible Shapes}

Which regular $n$-gons can be constructed?

\begin{definition}[Constructible]
    $a$ is constructible if you can construct a line segment of length $a$.
\end{definition}

\begin{definition}[Constructible Point]
    A point in $\rr^2$ is constructible if its coordinates are constructible.
\end{definition}

\begin{definition}[Constructible Line]
    A constructible line is a line determined by constructible points.
\end{definition}

\begin{theorem}{
    Constructible numbers lie in field extensions of $\q$.
}
    Suppose $a,b$ are constructible.
    Then they are closed under addition, subtraction, multiplication, and division by nonzero elements.
    Hence constructible numbers lie in a field extension of $\q$.
\end{theorem}

\begin{theorem}{
    If $a$ is constructible and $a\geq 0$, then $\sqrt{a}$ is constructible.
}
    Suppose a triangle is enclosed in a semicircle with triangle side length $1$ and diameter $\frac{a+1}{2}$.
    The relevant distance $x$ satisfies
    \[
        x^2=\left(\frac{a+1}{2}\right)^2-\left(\frac{a-1}{2}\right)^2.
    \]
    Thus
    \[
        x^2=a.
    \]
    So the distance $x=\sqrt{a}$ is constructible.
\end{theorem}

Suppose we have constructible points.
How do we get new points by intersecting lines, circles, and lines with circles?

\begin{definition}[New Constructible Points]
    If $\alpha$ is obtained from a constructible configuration, then adjoining $\alpha$ gives a quadratic extension.
    That is, any constructible point lies in a tower of fields
    \[
        \Q\subseteq\Q[a_1]\subseteq\Q[a_1,a_2]\subseteq\hdots\subseteq\ff,
    \]
    where
    \[
        \ff_k=\ff_{k-1}[a_k]
    \]
    and
    \[
        [\ff_k:\ff_{k-1}]=2.
    \]
\end{definition}

Let $\alpha$ be the root of a quadratic polynomial.
So every constructible number lies in a field $\ff$ where
\[
    [\ff:\Q]=2^n
\]
for some $n$.
Therefore $\sqrt[3]{2}$ is not constructible.

\begin{lemma*}[Constructible Points]{
    Let $r\in\mathbb{R}$.
    Then $r$ is constructible with straightedge and compass if and only if $r$ lies in a field extension $\mathbb{E}$ with
    \[
        [\mathbb{E}:\mathbb{Q}]=2^n
    \]
    for some $n$.
}
\end{lemma*}

$\pi$ is not constructible, and neither is it algebraic.
Therefore constructible points are only possible if
\[
    [\mathbb{Q}(r):\mathbb{Q}]=2^k
\]
for some $k$.
We will show that we cannot trisect $60^\circ$, since we can construct $60^\circ$, implying that not every angle can be trisected.
Because $20^\circ=\theta=\frac{\pi}{9}$ would need to be constructible, the number $\cos\theta$ would need to be constructible.
\begin{align*}
    \cos 2\theta &= \cos^2\theta-\sin^2\theta\\
    &= 2\cos^2\theta-1,
\end{align*}
and
\[
    \cos 3\theta=4\cos^3\theta-3\cos\theta.
\]
If $\theta=\frac{\pi}{9}$, then
\[
    \cos 3\theta=\cos\frac{\pi}{3}=\frac{1}{2}.
\]
Let $x=\cos 20^\circ$.
Then
\begin{align*}
    \frac{1}{2} &= 4x^3-3x,\\
    0 &= 8x^3-6x-1.
\end{align*}
We claim that $8x^3-6x-1$ is irreducible over $\mathbb{Q}$.
By the rational root test, the only possible rational roots are
\[
    \pm 1,\ \pm \frac{1}{2},\ \pm \frac{1}{4},\ \pm \frac{1}{8},
\]
and none of these are roots.
So the polynomial is irreducible over $\Q$.
Hence
\[
    [\Q(x):\Q]=3,
\]
which is not a power of $2$.
Therefore $x=\cos 20^\circ$ is not constructible, so $60^\circ$ cannot be trisected.Question:
Which regular $n$-gons can be constructed?
That is, for which $n$ can the angle
\[
    \frac{2\pi}{n}
\]
be constructed.
Such an angle can be constructed if and only if
\[
    \cos \frac{2\pi}{n}
    \quad\text{and}\quad
    \sin \frac{2\pi}{n}
\]
can be constructed.
\[
    \definecolor{zzttqq}{rgb}{0.6,0.2,0}
    \definecolor{ududff}{rgb}{0.30196078431372547,0.30196078431372547,1}
    \begin{tikzpicture}[line cap=round,line join=round,>=triangle 45,x=1cm,y=1cm]
\clip(-3.1902986457244733,-2.936287863782368) rectangle (3.569554280433086,2.160548852751021);
\fill[line width=2pt,color=zzttqq,fill=zzttqq,fill opacity=0.10000000149011612] (-1,-1) -- (1,1) -- (1,-1) -- cycle;
\draw [line width=2pt,color=zzttqq] (-1,-1)-- (1,1);
\draw [line width=2pt,color=zzttqq] (1,1)-- (1,-1);
\draw [line width=2pt,color=zzttqq] (1,-1)-- (-1,-1);
\draw (-0.2723211297635221,-0.8883141982266557) node[anchor=north west] {$\cos\frac{2\pi}{n}$};
\draw (1.0596316677595505,0.06638029248352595) node[anchor=north west] {$\sin\frac{2\pi}{n}$};
\draw (-0.2954185771194135,0.30505391516107133) node[anchor=north west] {1};
\draw (-0.626481989220524,-0.5033567422951308) node[anchor=north west] {$\frac{2\pi}{n}$};
\draw (0.6,-0.4033567422951308) node[anchor=north west] {$\frac{\pi}{2}$};
\begin{scriptsize}
    \draw [fill=ududff] (-1,-1) circle (2.5pt);
    \draw[color=ududff] (-0.9421471030843737,-0.8382697289555574) node {$A$};
    \draw [fill=ududff] (1,1) circle (2.5pt);
    \draw[color=ududff] (1.0596316677595505,1.1635090418883718) node {$B$};
    \draw[color=zzttqq] (0.11263632616800179,0.008636674093797217) node {$c$};
    \draw[color=zzttqq] (0.905648685386941,0.09332731439873268) node {$a$};
    \draw[color=zzttqq] (0.02794568586306654,-0.784375685125144) node {$b$};
\end{scriptsize}
    \end{tikzpicture}
\]
Thus a regular $n$-gon is constructible if and only if
\[
    \left(\cos \frac{2\pi}{n},\sin \frac{2\pi}{n}\right)
\]
is a constructible point.
If we let
\[
    \rho=\cos \frac{2\pi}{n}+i\sin \frac{2\pi}{n},
\]
then
\[
    \rho^n=1.
\]
So $\rho$ is an $n$th root of unity satisfying
\[
    x^n-1=0.
\]
Suppose
\[
    n=2^{a_1}p_2^{a_2}\cdots p_k^{a_k}
\]
is a prime factorization, where $p_1=2$ and $p_j$ is odd for $j\geq 2$.
Then the regular $n$-gon is constructible if and only if each odd prime factor occurs to the first power and each such odd prime is a Fermat prime.

\begin{definition}[Fermat Prime]
    If
    \[
        2^{2^k}+1=p
    \]
    is prime, then $p$ is a Fermat prime.
\end{definition}

\begin{corollary}{
    The constructible regular $n$-gons are exactly those for which
    \[
        n=2^m p_{1}p_2\cdots p_r,
    \]
    where the $p_i$ are distinct Fermat primes.
}
\end{corollary}

\begin{proposition}{
    Let $\phi:\mathbb{Q}[x]\to \mathbb{Q}[\alpha]$ be defined by
    \[
        \phi(p(x))=p(\alpha).
    \]
    Then $\phi$ is surjective.
}
    Suppose $\beta\in\mathbb{Q}[\alpha]$.
    Then by definition of $\mathbb{Q}[\alpha]$, there exists $f(x)\in\mathbb{Q}[x]$ such that
    \[
        \beta=f(\alpha).
    \]
    Thus
    \[
        \phi(f(x))=f(\alpha)=\beta.
    \]
    So $\phi$ is surjective.
\end{proposition}

$\phi$ is injective if and only if $\ker\phi=\{0\}$.
This is a consequence of the first isomorphism theorem.

\begin{proposition*}{
    \[
        \ker\phi=\{0\}
        \iff
        (\forall f(x)\in\Q[x],\ f(\alpha)=0 \Rightarrow f(x)=0)
        \iff
        \alpha/\mathbb{Q}\text{ is transcendental.}
    \]
}
\end{proposition*}